% Figure arealaw2d_tilt_overlap.
% Owner: none (note-only figure).
% Used by: gapped-2d-notes AL Sec. IV A (Lemma 4, concentration of Schmidt spectra),
% Eq. al:s3:overlap.
% Source: arXiv OpenAI2026AreaLaw, sections/02-initial.tex:60-93 (Theta, F) and
% 174-185 (the overlap M(u/2)^2/M(u) = |<Psi,Theta>|^2).
% Formula: <Psi, F Psi> = Tr(rho rho^{-u/2}) = Tr rho^{1-u/2}, where
% rho = Tr_{B^c}|Psi><Psi| and F = rho^{-u/2} on B. Powers are taken on the
% support of rho; F is the identity on ker rho, where Psi has no weight, so
% rho F = rho^{1-u/2} on the whole space. With ||F Psi||^2 = Tr rho^{1-u} and
% Theta = F Psi / ||F Psi|| this gives
% |<Psi,Theta>|^2 = (Tr rho^{1-u/2})^2 / Tr rho^{1-u}.
% Ink-to-index: panel 1 has two wires, B^c (top) and B (bottom), named west of
% the bra. The traced wire is on top, as in polypeps_trace_stripping, so that F
% and every operator name sit south (one station) with no wire below them. The
% ket Psi (east) and the bra Psi^dagger (west) span both wires, names south as
% in tn2d_single_conjugation; the top wire joins them directly (the partial
% trace that makes rho), and the bottom wire carries F = rho^{-u/2}. Panel 2 is
% a trace on H_B: the arc is the trace, and the boxes read left to right as the
% product inside Tr. The arc leaves the boxes' own faces (no extenders) so that
% the display fits one column.
% Joiners: the first = is the partial trace plus cyclicity (cyc, no hypothesis);
% the second is functional calculus on the support of rho, written as the
% statement proved (J4).
% Boundary signature: empty (scalars).
% Fit: C (\tnshow[6mm]); one line in the PRL column at 6 mm.
% Label macros: none.
% Related: polypeps_trace_stripping (the same partial-trace-to-ring first move,
% a different identity: conjugation, not a power).
\begin{tenkzequation}
    \tenkzkernel
    \begin{tenkz}[rows={wire,wire}, cols=5, bonds=none, size=l]
        \tn[at=(1,1), name=toBra, skin=box, wires=2, label pos=s,
            ports={0@1:physical, 0@2:physical}]{\bra{\Psi}}
        \tn[at=(2,3), name=toF, skin=box, label pos=s,
            ports={180:physical, 0:physical}]{\rho^{-u/2}}
        \tn[at=(1,5), name=toKet, skin=box, wires=2, label pos=s,
            ports={180@1:physical, 180@2:physical}]{\ket{\Psi}}
        \tnwire{toBra.0@1}{toKet.180@1}
        \tnwire{toBra.0@2}{toF.180}
        \tnwire{toF.0}{toKet.180@2}
        \tnmark[form=label, label pos=w]{(1,1)}{$B^c$}
        \tnmark[form=label, label pos=w]{(2,1)}{$B$}
    \end{tenkz}
    $\overset{\text{cyc}}{=}$
    \begin{tenkz}[rows={wire}, cols=3, bonds=none, size=l]
        \tn[at=(1,1), name=toRho, skin=box, label pos=s,
            ports={180:physical, 0:physical}]{\rho\vphantom{\rho^{-u/2}}}
        \tn[at=(1,3), name=toFF, skin=box, label pos=s,
            ports={180:physical, 0:physical}]{\rho^{-u/2}}
        \tnwire{toRho.0}{toFF.180}
        \tnwire[route=arc]{toFF.0}{toRho.180}
    \end{tenkz}
    $=\operatorname{Tr}\rho^{1-u/2}$
\end{tenkzequation}
